\documentclass[10pt,a4paper]{amsart}
\usepackage[margin=19mm]{geometry}
\usepackage{amsmath,amssymb,mathtools}
\usepackage[scr=rsfs,cal=euler]{mathalfa}
\usepackage{xfrac}
\usepackage[unicode,hidelinks,bookmarks=false]{hyperref}
\newcommand{\ie}{i.e.\ }
\newcommand{\cf}{cf.\ }
\newcommand{\resp}{respectively\ }
\newcommand{\Subsection}{\subsection}
\providecommand{\nobreakdash}{\nobreak\hbox{-}}
\setlength{\emergencystretch}{2em}
\multlinegap0pt
\usepackage{tikz}
\usetikzlibrary{cd}
\usepackage{mathtools}
\usepackage{xcolor}
\usepackage{amssymb}
\usepackage[shortlabels]{enumitem}
\newtheorem{definition}{Definition}
\newtheorem{prop}{Proposition}
\newtheorem{hypothesis}{Hypothesis}
\newtheorem{theorem}{Theorem}
\newtheorem{remark}{Remark}
\newtheorem{lemma}{Lemma}
\newtheorem{cor}{Corollary}
\DeclareMathOperator{\Boxe}{Box}
\newcommand{\G}{\mathbf{G}}
\DeclareMathOperator{\Hom}{Hom}
\newcommand{\K}{\mathbf{K}}
\newcommand{\Le}{\mathbf{L}}
\DeclareMathOperator{\NS}{NS}
\DeclareMathOperator{\Pic}{Pic}
\newcommand{\Q}{\mathbf{Q}}
\newcommand{\Z}{\mathbf{Z}}
\DeclareMathOperator{\age}{age}
\DeclareMathOperator{\ext}{ext}
\DeclareMathOperator{\grp}{grp}
\DeclareMathOperator{\num}{num}
\DeclareMathOperator{\orb}{orb}
\newcommand{\sector}{\pi_0(\sheafI_{\mu} X)}
\newcommand{\sheafI}{\mathcal{I}}
\newcommand{\sheafO}{\mathcal{O}}
\newcommand{\stackS}{\mathcal{S}}
\newcommand{\stackT}{\mathcal{T}}
\newcommand{\stackY}{\mathcal{Y}}
\newcommand{\stackZ}{\mathcal{Z}}
\newcommand{\twistsector}{\pi_0^*(\sheafI_{\mu} X)}
\title{full title}
\author{Nicolas Bongiorno}
\date{Source: arXiv:2510.22325v3 (CC BY 4.0)}
\begin{document}
\maketitle

\noindent\emph{Source and licence.} full title. Version arXiv:2510.22325v3 (CC BY 4.0), distributed by arXiv under the Creative Commons Attribution 4.0 International licence, http://creativecommons.org/licenses/by/4.0/, as recorded in the arXiv licence metadata for this exact version and shown on the abstract page https://arxiv.org/abs/2510.22325v3. Full licence notice: \texttt{licenses/arxiv-2510.22325-CC-BY-4.0.txt}.\par\smallskip

\noindent\textit{Editorial scope.} Counted unit: the article's prop prop\_surjectivite\_local\_extended\_quotient\_description (source0.tex lines 2913--2920). The article's complete original proof of it is delivered with it (source0.tex lines 2923--3072, one \texttt{proof} environment of 150 lines). No mathematics is written, repaired or re-derived: the statement and the proof are the article's own lines, in the article's own notation, and no retained line is substituted or re-flowed.  Retained uncounted as the source-local context the proof needs: lines 1239--1241; lines 1267--1270; lines 1288--1291; lines 1478--1487; lines 1588--1593; lines 2255--2258; lines 2312--2318; lines 2364--2379; lines 2383--2392; lines 2412--2420; lines 2696--2698; lines 2703--2718; lines 2889--2894; lines 5536--5538.  Nothing was omitted from any retained span, so the package carries no omission record.  No retained line contains a citation, so the wrapper carries no bibliography and no bibliography entry.  The retained lines contain the article's own diagram code, which the wrapper typesets with the standard tikz packages; no image file is delivered and the package declares no asset dependency.  Editorial formatting only, in addition: an AMS \texttt{amsart} preamble that restates the article's own declarations and macros used by the retained lines, namely the environments \texttt{definition}, \texttt{prop}, \texttt{hypothesis}, \texttt{theorem}, \texttt{remark}, \texttt{lemma}, \texttt{cor} on separate counters, together with the standard packages those lines need.  The retained environments print in this document as Definition 1, Proposition 1, Hypothesis 1, Theorem 1, Remark 1, Lemma 1, Corollary 1 in order of appearance, whereas the article prints the same items with the numbering of its own full document.  The complete native source of the article as distributed in the arXiv e-print for this exact version is bundled unchanged as \texttt{sources/187783/source0.tex}.
\begin{definition}\label{def_age_univ_torsor}
   For $y \in \stackT(\K)$, we define $$\age_T(y,-) \in \left(X^*(T)\right)^{\vee}_{\Q} = \Hom_{\grp}(X^*(T),\Q)$$ such that for any $\chi \in X^*(T)$, $\age_T(y,\chi)$ is the $\Q$-value $\frac{v_{\K'}(\chi(t))}{e(\K' / \K)} $ where $$z = t.y \in \stackT(\sheafO_{\K'})$$ lies in the $T$-orbit of $y$ for some finite extension $\K' / \K$. We shall denote this by $\age_{v,T}(y,-)$ when we want to emphasize the valuation $v$.
\end{definition}

\begin{prop}\label{prop_age_and_logarithm}
    Let $y \in \stackT(\K)$ and $t \in T(\K)$. We have:
    $$\age_v(t.y,-) = \age_v(y,-) - \log_{T,v}(t) .$$
\end{prop}

\begin{hypothesis}\label{hypothesis_zero_of_local_degree}
    For any finite field extension $\K' / \K$, we have:
    $$\stackT(\sheafO_{\K'}) = \{y \in \stackT(\K') \mid \age_{v_{\K'}}(y,-) = 0 \} . $$
\end{hypothesis}

\begin{definition}\label{definition_general_age_pairing}
For any character $\chi \in X^*(T)$, we define
\[
\age_{\alpha,T}(b,\chi) = b^*\bigl(\alpha^*(\chi)\bigr) \in \tfrac{1}{N}\Z/\Z.
\]
This defines an age pairing associated with the multiplicative group $T$,
\[
\age_{\alpha,T} : \sector \times X^*(T) \longrightarrow \Q/\Z.
\]
\end{definition}

\begin{theorem}\label{theorem_age_extended_universal_torsor}
For any $\chi \in X^{*}(T)$, we have:
\[
\age_{\alpha,T}\bigl(\psi_v(P),\chi\bigr) \equiv \age_{v,T}(y,\chi) \mod \Z .
\]
\end{theorem}

\begin{prop}\label{prop_age_divisor_associated_to_an_edge}
For $\sigma \in \Sigma_{\max}$ and for any $b \in \Boxe(\sigma) \subset \sector$ such that $$\overline{b} = \sum\limits_{\rho \in \sigma(1)} q_{\rho} . \overline{b_{\rho}} \in N_{\Q},$$ we have for any $\rho \in \sigma(1)$:
$$\age_{\num}(\stackY_b,-[D_{\rho}]) = q_{\rho} .$$Moreover, if $\rho \not\in \sigma(1)$, $\age(\stackY_b,[D_{\rho}]) = 0 \mod \Z$.
\end{prop}

\begin{remark}\label{remark_exact_sequence_dual_orbifold_picard_group}
    In particular, in the case where $\mathcal{C} = \twistsector$, we shall write $$\beta_{\twistsector} = \beta^{\ext} .$$Since $\beta^{\ext}$ is surjective, we obtain the exact sequence
\[
0 \longrightarrow \Pic_{\orb}(X)^{\vee} \xrightarrow{\pi^{\vee}} \Z^{\Sigma(1)} \oplus \Z^{\twistsector} \xrightarrow{\beta^{\ext}} N \longrightarrow 0 \, .
\]This also proves that $\Pic_{\orb}(X)$ is torsion free (see Corollary \ref{corollaire_orbifold_picard_group_torsion_free} for an other proof). 

\end{remark}

\begin{definition}\label{definition_extended_universal_torsor}
The quotient morphism
\[
q^{\ext} : \stackT \times \G_m^{\twistsector} \longrightarrow X,
\]
arising from the quotient presentation of $X$, where $T_{\NS,\orb}$ acts on
$\stackT \times \G_m^{\twistsector}$ via the inverse of the action defined at the
beginning of this paragraph, defines a $T_{\NS,\orb}$-torsor. We call it the
\emph{extended universal torsor}.

For $\stackS \in \sector$ and $(L,\varphi) \in \Pic_{\orb}(X)$, we denote by
\[
\age\bigl(\stackS,(L,\varphi)\bigr)
\]
the age pairing defined via this torsor (see Definition~\ref{definition_general_age_pairing}).
\end{definition}

\begin{prop}\label{proposition_functoriality_age}
The type of the extended universal torsor is given by the natural morphism
\[
(L,\varphi) \longmapsto L.
\]
Moreover, for any $\stackS \in \sector$ and any $(L,\varphi) \in \Pic_{\orb}(X)$, we have
\[
\age\bigl(\stackS,(L,\varphi)\bigr) = \age(\stackS,L).
\]
\end{prop}

\begin{prop}\label{prop_extended_univ_torsor_zero_age}
    Let $Y \rightarrow X$ a $T$-torsor other the toric stack $X$, which is either:
    \begin{itemize}
        \item the $T_{\NS}$-torsor given by $Y = \stackT_{\Sigma}$ when $\Pic(X)$ is torsion free;
        \item the $T_{\NS,\orb}$-torsor given by $Y = \stackT_{\Sigma} \times \G_m^{\twistsector}$.
    \end{itemize}
     Then hypothesis \ref{hypothesis_zero_of_local_degree} is verified: for any finite field extension $\K' / \K$, we have:
    $$\stackT(\sheafO_{\K'}) = \{y \in \stackT(\K') \mid \age_{v_{\K'}}(y,-) = 0 \} . $$
\end{prop}

\begin{equation}\label{equation_formula_equivariance_valuation}
v_{\Le}(t \cdot y) = v_{\Le}(y) - \pi^{\vee}\left(\log_{v_{\Le}}(t)\right)    
\end{equation}by Definition \ref{definition_extended_universal_torsor}.

\begin{prop}\label{proposition_computation_residue_map_extended_parametrisation}
Let $y \in \stackT^{\ext}(\Z_v)$ and set $P = q^{\ext}(y)$. Then, for $\stackS \in \twistsector$, we have the equivalence
\[
p_v \mid y_{\stackS}
\;\Longleftrightarrow\;
\psi_v(P) = \stackS \, .
\]
In particular, for $y \in \stackT^{\ext}(\Z_v)$ such that $p_v \mid y_{\stackS}$, we have
\[
\pi^{\vee}\!\left( \age_v(y,-) \right)
= e_{\stackS}
-
\sum_{\rho \in \sigma_\stackS(1)}
\age_{\num}\!\left(\stackS,-[D_{\rho}]\right)\, e_{\rho} \, .
\]
\end{prop}

\begin{lemma}\label{lemma_inequality_valuation_plus_age}
Let $y \in \stackT^{\ext}(\Q_v)$. Then for every $\rho \in \Sigma(1)$, we have the inequality
\[
v(y_{\rho}) + \lfloor \age_v(y,-[D_{\rho}]) \rfloor \geqslant 0 .
\]
\end{lemma}

\begin{cor}\label{corollaire_orbifold_picard_group_torsion_free}
    $\Pic_{\orb}(X)$ is torsion free.
\end{cor}

\begin{prop}\label{prop_surjectivite_local_extended_quotient_description}
Let $y \in \stackT^{\ext}(\Q_v)$. Then there exists an element
$\alpha \in \Pic_{\orb}(X)^{\vee}$ such that, if
$t \in T_{\NS,\orb}(\Q_v)$ satisfies $\log_v(t) = \alpha$, then
\[
t \cdot y \in \stackT^{\ext}(\Z_v).
\]
\end{prop}

\begin{proof}
Let $P = q^{\ext}(y)$. Assume first that $\psi_v(P) = 0$. In this case, by Theorem~\ref{theorem_age_extended_universal_torsor}, we have
\[
\age_v(y,-) \in \Pic_{\orb}(X)^{\vee},
\]
so we may consider an element $t \in T_{\NS,\orb}(\Q_v)$ such that
\[
\log_v(t) = \age_v(y,-).
\]
Then, by Proposition~\ref{prop_age_and_logarithm}, we obtain
\[
\age_v(t \cdot y,-) = 0.
\]
It therefore follows from Proposition~\ref{prop_extended_univ_torsor_zero_age} that
\[
t \cdot y \in \stackT(\Z_v) \times \G_m(\Z_v)^{\twistsector} \subset \stackT^{\ext}(\Z_v).
\]

Assume now that $\psi_v(P) = \stackS \in \twistsector$. Let $\Le/\Q_v$ be a finite extension such that there exists
\[
z \in \stackT(\sheafO_{\Le}) \times \G_m(\sheafO_{\Le})^{\twistsector}
\]
lying in the same $T_{\NS,\orb}(\Le)$-orbit as $y_{|\Le}$. Using the definition of the lift of the age associated with the $T_{\NS,\orb}$-torsor (see Definition~\ref{def_age_univ_torsor}), we obtain the equality
\[
\frac{v_{\Le}(z)}{e(\Le/\Q_v)}
=
v(y) - \pi^{\vee}(\age_v(y,-)) 
\]
in $\left( \Z \cup \{+\infty\} \right)^{\Sigma_{\ext}(1)}$.

By
Theorem~\ref{theorem_age_extended_universal_torsor}, Proposition \ref{proposition_functoriality_age} and
Proposition~\ref{prop_age_divisor_associated_to_an_edge}, we further have
\begin{equation}\label{equation_computation_value_n}
\pi^{\vee}\!\left( \age_v(y,-) \right)
=
n - \sum_{\rho \in \sigma_\stackS(1)} q_{\rho}\, e_{\rho} \, ,   
\end{equation}where
$q_{\rho} = \age_{\num}\!\left(\stackS,-[D_{\rho}]\right) \in [0,1[$ and
$n \in \Z^{\Sigma_{\ext}(1)}$.

We denote by $b_{\stackS} \in \Boxe(\Sigma)$ the element associated with $\stackS$.
Recall that
\[
\overline b_{\stackS}
=
\sum_{\rho \in \sigma_\stackS(1)} q_{\rho}\, \beta^{\ext}_{\Q}(e_{\rho}).
\]
By Remark~\ref{remark_exact_sequence_dual_orbifold_picard_group}, we have the following
commutative diagram with exact rows:
\begin{center}
\begin{tikzcd}
0 \arrow[r]
& \Pic_{\orb}(X)^{\vee}
\arrow[r, "\pi^{\vee}"]
\arrow[d]
& \Z^{\Sigma_{\ext}(1)}
\arrow[r, "\beta^{\ext}"]
\arrow[d]
& N
\arrow[r]
\arrow[d]
& 0
\\
0 \arrow[r]
& \Pic_{\orb}(X)^{\vee}_{\Q}
\arrow[r]
& \Q^{\Sigma_{\ext}(1)}
\arrow[r]
& N_{\Q}
\arrow[r]
& 0
\end{tikzcd}
\end{center}

From the previous equality, we obtain
\[
\beta^{\ext}_{\Q}(n) = \overline b_{\stackS}.
\]
Therefore, there exists $\stackS' \in \twistsector$ such that
$\sigma_\stackS' = \sigma_\stackS$,
$\overline b_{\stackS'} = \overline b_{\stackS}$, and
\[
\beta^{\ext}(n) = b_{\stackS'}.
\]
It follows that $n - e_{\stackS'} \in \pi^{\vee}\left(\Pic_{\orb}(X)^{\vee}\right)$, thus there exists $\alpha \in \Pic_{\orb}(X)^{\vee}$ such that
\[
n - e_{\stackS'} = \pi^{\vee}\left(\alpha\right).
\]

We then consider $t \in T_{\NS,\orb}(\Q_v)$ such that
\[
\log_v(t) = \alpha.
\]
Using \eqref{equation_formula_equivariance_valuation}, we obtain the equality
\[
\frac{v_{\Le}(z)}{e(\Le/\Q_v)}
=
v(t \cdot y)
+
\sum_{\rho \in \sigma_\stackS(1)} q_{\rho}\, e_{\rho}
-
e_{\stackS'}
\]
in $\left( \Z \cup \{+\infty\} \right)^{\Sigma_{\ext}(1)}$. We now show that $y' = t \cdot y \in \stackT^{\ext}(\Z_v)$. First, we observe that for
every $\stackZ \neq \stackS'$, we have $v(y'_{\stackZ}) = 0$, and that
$v(y'_{\stackS'}) = 1$. It therefore remains to prove that for any primitive
collection $C$, we have
\[
\min_{\rho \in C \setminus \sigma_{\stackS'}(1)} v(y'_{\rho}) = 0 .
\]

To begin with, since for every $\rho \in \Sigma(1)$ we have by \eqref{equation_computation_value_n}
\[
n(e_{\rho}) = - \lfloor \age_v(y,-[D_{\rho}]) \rfloor,
\]
it follows that
\[
v(y'_{\rho})
=
v(y_{\rho}) + \lfloor \age_v(y,-[D_{\rho}]) \rfloor
\geqslant 0
\]
by Lemma~\ref{lemma_inequality_valuation_plus_age}. Now, using the fact that
\[
z \in \stackT(\sheafO_{\Le}) \times \G_m(\sheafO_{\Le})^{\twistsector},
\]
we have
\[
0
=
\min_{\rho \in C}
\frac{v_{\Le}(z_{\rho})}{e(\Le/\Q_v)}
=
\min_{\rho \in C}
\left(
v(y'_{\rho}) + \age_{\num}\!\left(\stackS',-[D_{\rho}]\right)
\right).
\]
Since for every $\rho \notin \sigma_{\stackS'}(1)$ we have
$\age_{\num}\!\left(\stackS',-[D_{\rho}]\right) = 0$, while for
$\rho \in \sigma_{\stackS'}(1)$ we have
$\age_{\num}\!\left(\stackS',-[D_{\rho}]\right) \in ]0,1[$, and since moreover
$v(y'_{\rho}) \geqslant 0$, it follows that
\[
\min_{\rho \in C \setminus \sigma_{\stackS'}(1)} v(y'_{\rho}) = 0 .
\]Note that, by Proposition~\ref{proposition_computation_residue_map_extended_parametrisation},
we obtain $\stackS' = \stackS$.

\end{proof}

\end{document}
